\section{三个实证发现}
\label{sec:findings}

下文所有数字都可由开放脚本 \texttt{heal\_baseline.py}、\texttt{false\_heal\_probe.py} 与 \texttt{cross\_app\_probe.py} 在 \texttt{healreact/bench/cases/} 下的 JSON 工件上复现。本地 Qwen 批量实验（两个 L3 变体 + 两个误修复探针）在 Apple M5 Pro / 48\,GB 上仅用 Apple Metal 加速的实际物理耗时约为 8\,min；\texttt{gpt-4o} 对照经 OpenAI 兼容接口单独运行。

\subsection{复现配置}
为便于复现，本文固定以下设置：L3 修复器与意图标注器分别经 Ollama 运行 \texttt{qwen2.5-coder:7b} 与 \texttt{qwen2.5:3b}，二者均取 \texttt{temperature}=0（贪心解码、单次采样）；候选检索采用加性 token 重叠打分，权重为 testId$\times 3$、dataAttr key/value$\times 2$、aria/id/class/text/path$\times 1$，取 top-$k$ 候选（$k$ 与各探针配置一并记录在工件 JSON 中）。为应对模型单一带来的有效性威胁，修订版还在同一 F1 协议上加入一个前沿闭源模型对照（\texttt{gpt-4o}，OpenAI 兼容 \texttt{/chat/completions} 接口，\texttt{temperature}=0，单次采样；正式工件使用 2\,s 请求间隔与 429 退避重试，最终错误数为 0）。每个失效用例的逐条记录（用例 id、koenig 所属 commit、断裂选择器、ReproBreak \texttt{new\_locator}、模型给出的选择器与所选候选下标、生成理由、选择器语法校验结果、解析得到的目标元素、是否判为误修复/精确匹配）以及汇总头部的模型标识均随 \texttt{healreact/bench/cases/} 工件一并发布，作为本文结果的复现与审计依据。\textbf{代码与数据可用性}：上述脚本、JSON 工件、工件清单与论文快照已公开归档于 GitHub: \url{https://github.com/vgawen/healreact}；表~\ref{tab:falseheal}--\ref{tab:cross-app} 中的每一项原始计数（样本数、放弃/误修复/未解析、精确匹配）均可由该归档直接复算。工件清单汇总各脚本的运行命令、输入数据校验和、koenig 与 payload 的 commit 标识及模型版本。

\subsection{共享 L3 协议}
在 F1 与 F2 中，L3 修复器默认经由 Ollama 以 \texttt{temperature}=0 调用 \texttt{qwen2.5-coder:7b}；F1 的模型泛化对照则在不改变候选、提示格式和解析器事实标准的前提下替换为 \texttt{gpt-4o}。用户提示包含: (i) 原样给出的断裂 Playwright 选择器；(ii) 失效行附近的 $\pm 3$ 行测试源码；(iii) 来自 LocatorSheet 的 top-$k$ 候选，按加性 token 重叠打分检索（testId $\times 3$、dataAttr key/value $\times 2$、aria/id/class/text/path $\times 1$）。模型必须给出三行: \texttt{SELECTOR}、\texttt{CANDIDATE\_IDX}、\texttt{RATIONALE}（纯文本，非 JSON: 在模型反复在选择器中输出未转义引号后放弃了 JSON）。解析器的多匹配冲突处理规则为：按 \texttt{(componentFile, line)} 顺序的第一个匹配记录作为解析目标；多匹配的选择器会被标记，但指标计算使用第一个命中。

\subsection{F1: 目标缺失下的隐性误修复}
\label{sec:f1}

\subsubsection{设置}
对 75 个\emph{可达}的 koenig 失效用例（即失效 commit 处 LocatorSheet 中存在真值（ground truth）元素的用例），本文以\emph{相同}的提示与检索条件重跑 L3 修复器，但将真值元素从 LocatorSheet 与候选集中\emph{强制移除}，使模型在字面意义上无法选中正确答案。此时正确响应应当是放弃（\texttt{CANDIDATE\_IDX: -1}）；任何给出可解析选择器的响应都是一次会沉默地将测试改写到错误元素上的隐性误修复。

我们比较两个提示变体:
\begin{enumerate}
\item \textsc{Vanilla}：与 L3 基线相同的系统提示。
\item \textsc{Abstain}: 同一提示外加一条明确的“如无候选契合则放弃”守门规则，告诉模型 \emph{“一次沉默地把测试改写到错误元素上的盲修复永远比放弃更糟”}。
\end{enumerate}

\subsubsection{结果}
表~\ref{tab:falseheal} 给出结果。

\begin{table}[t]
\centering
\caption{误修复探针（\S\ref{sec:f1}）。所有模型均使用相同的 75 个可达 koenig 用例与相同的 ground-truth 移除协议。}
\label{tab:falseheal}
\small
\resizebox{\columnwidth}{!}{%
\begin{tabular}{llrrrr}
\toprule
模型 & prompt & 放弃\,(好) & 误修复\,(坏) & 未解析 & \textbf{误修复\,\%} \\
\midrule
\texttt{qwen2.5-coder:7b} & \textsc{Vanilla} & 13 & 59 & 3 & \textbf{78.7\%} \\
\texttt{qwen2.5-coder:7b} & \textsc{+ 放弃守门规则} & 13 & 59 & 3 & \textbf{78.7\%} \\
\texttt{gpt-4o} & \textsc{Vanilla} & 12 & 63 & 0 & \textbf{84.0\%} \\
\texttt{gpt-4o} & \textsc{+ 放弃守门规则} & 24 & 51 & 0 & \textbf{68.0\%} \\
\bottomrule
\end{tabular}}
\\[2pt]
{\scriptsize $n=75$ 个可达 koenig 失效用例。Ground truth 目标已从候选集中移除；任何被给出的可解析选择器即为一次隐性误修复。Qwen 的 78.7\%（59/75）Wilson 95\% 置信区间约为 [68.1\%, 86.4\%]；GPT-4o Abstain 的 68.0\%（51/75）约为 [56.8\%, 77.5\%]。}
\end{table}

为排除“聚合计数偶然相同”的可能，表~\ref{tab:falseheal-paired} 给出 Qwen 两提示变体的\emph{逐用例}配对交叉表。

\begin{table}[t]
\centering
\caption{Vanilla 与 Abstain 的逐用例配对交叉表（\S\ref{sec:f1}，$n=75$）。所有非对角元素均为 0，表明两变体在每一个用例上给出\emph{完全相同}的判定，而非仅聚合计数相同。}
\label{tab:falseheal-paired}
\small
\begin{tabular}{lrrr}
\toprule
Vanilla＼Abstain & 放弃 & 误修复 & 未解析 \\
\midrule
放弃\,(好)   & 13 & 0 & 0 \\
误修复\,(坏) & 0  & 59 & 0 \\
未解析       & 0  & 0 & 3 \\
\bottomrule
\end{tabular}
\\[2pt]
{\scriptsize 在 \texttt{temperature}=0 的确定性解码下，75 个用例的逐例标签两两一致（off-diagonal $=0$，$\Delta=0$）。}
\end{table}

\begin{table}[t]
\centering
\caption{\texttt{gpt-4o} Vanilla 与 Abstain 的逐用例配对交叉表（\S\ref{sec:f1}，$n=75$）。Abstain 提示把 12 个 Vanilla 误修复改为放弃，但仍有 51 个用例保持误修复。}
\label{tab:falseheal-gpt-paired}
\small
\begin{tabular}{lrrr}
\toprule
Vanilla＼Abstain & 放弃 & 误修复 & 未解析 \\
\midrule
放弃\,(好)   & 12 & 0 & 0 \\
误修复\,(坏) & 12 & 51 & 0 \\
未解析       & 0  & 0 & 0 \\
\bottomrule
\end{tabular}
\\[2pt]
{\scriptsize 与 Qwen 不同，\texttt{gpt-4o} 会响应放弃提示；但 51/75 的误修复残留说明软提示仍不足以成为无人值守 CI 的安全边界。}
\end{table}

\subsubsection{解读}
对 \texttt{qwen2.5-coder:7b} 而言，78.7\% 是该协议下隐性误修复的\emph{对抗性上界}，因为每个用例都是目标缺失构造。在混合工作负载下，只有 L1 sheet 缺失目标的用例（本试点中是 18/93 = 19\% 的 koenig 失效）面临此风险；朴素地相乘得到 $\approx 0.19 \times 0.79 = 15\%$，作为一个粗略的风险估计、而非测量到的部署率（对抗性探针的行为不必一对一地迁移到真正只在运行时可达的用例上）。即便作为风险估计，我们仍认为这已经高到不能允许 CI 无人值守地自动提交 LLM 提议的选择器补丁。

Qwen 的 \textsc{Vanilla} 与 \textsc{Abstain} 等同性是更具诊断意义的发现，且这一等同性是\emph{逐用例}成立的（表~\ref{tab:falseheal-paired}：75 个用例的配对判定无一发生改变）：一个 7B 规模的代码模型在该任务上不会因软提示式的守门规则而改变其可观察行为。GPT-4o 对照使结论更细化而非被推翻：更强模型确实能利用放弃提示，将 12 个误修复改为放弃（表~\ref{tab:falseheal-gpt-paired}），但在同一对抗性协议下仍有 51/75 的用例输出错误可解析选择器。因而本文的强结论不是“所有模型必然忽视放弃提示”，而是：在缺乏外部确定性验证器时，仅凭提示文本无法作为可审计的安全边界；即便前沿模型能降低风险，仍需确定性拒绝机制或行为重放验证器来证明补丁没有错指。

\subsection{F2：L3 基线与两个低开销确定性手段}
\label{sec:f2}

\subsubsection{设置}
F2 仅作为辅助工程经验报告。在同样的 75 个可达 koenig 失效用例上，我们比较两杠杆均关的 v0 与两杠杆均开的 v1，并用二因子析因分离 testId 加权检索与强锚点后置过滤的边际贡献；所有臂均使用 \texttt{qwen2.5-coder:7b}、同一系统提示、同一候选池与同一解析器 oracle。Ground truth 是 ReproBreak 中 \texttt{new\_locator} 字符串所匹配的目标元素（经解析器解析）。

\subsubsection{结果}
表~\ref{tab:l3-levers} 显示，两杠杆全开时 top-1 精确元素匹配从 55/75（73.3\%）提升到 58/75（77.3\%），全量 93 个失效用例上的静态修复代理指标从 59.1\% 提升到 62.4\%，且未引入回归。表~\ref{tab:l3-factorial} 进一步说明，该改进来自两个确定性杠杆的叠加：单开 testId 加权或强锚点后置过滤均为 +2/-0，两者全开为 +3/-0；但这些差异在 $n=75$ 下均未达到统计显著。

\begin{table}[t]
\centering
\caption{L3 v0 $\to$ v1 端到端对比（75 个可达 koenig 失效）。v1 = v0 + testId 加权检索 + 强锚点后置过滤；v0 为两杠杆均关的统一权重基线。}
\label{tab:l3-levers}
\small
\resizebox{\columnwidth}{!}{%
\begin{tabular}{lrr}
\toprule
指标 & v0（两杠杆关） & v1（两杠杆开） \\
\midrule
top-1 精确元素匹配 & 55 / 75 (73.3\%) & \textbf{58 / 75 (77.3\%)} \\
有效 Playwright 选择器 & 68 / 75 (90.7\%) & 68 / 75 (90.7\%) \\
静态修复代理指标（全量 93） & 55 / 93 (59.1\%) & \textbf{58 / 93 (62.4\%)} \\
\midrule
强锚点后置过滤触发 & 0 & 1 / 75（正确） \\
新增正确用例（vs v0） & --- & 3 \\
回归用例（vs v0） & --- & 0 \\
\bottomrule
\end{tabular}}
\end{table}

\begin{table}[t]
\centering
\caption{两杠杆二因子析因（top-1 精确元素匹配，$n=75$）。行=testId 加权，列=强锚点后置过滤。每个杠杆单独带来 +2、两者合计 +3、零回归；方向一致但在该样本量下未达统计显著。}
\label{tab:l3-factorial}
\small
\begin{tabular}{lrr}
\toprule
 & 过滤：关 & 过滤：开 \\
\midrule
加权：关 & 55 / 75 & 57 / 75 \\
加权：开 & 57 / 75 & \textbf{58 / 75} \\
\bottomrule
\end{tabular}
\\[2pt]
{\scriptsize McNemar 精确检验（双侧）：单开加权或单开过滤相对“两关”均为 $+2/-0$（$p=0.50$）；两者全开相对“两关”为 $+3/-0$（$p=0.25$）。所有差异方向一致、零回归，但在 $n=75$ 下均未达 $p<0.05$。}
\end{table}

\subsubsection{解读}
两项手段均开销较低、完全确定性、无需任何模型再训练。借助表~\ref{tab:l3-factorial} 的二因子析因可分离两者贡献：以“两杠杆均关”为基线（55/75），单开 testId 加权检索与单开强锚点过滤各带来 +2 个精确匹配，两者全开合计 +3、且\emph{零回归}。须如实说明：这些增益方向一致、无一造成回归，但在 $n=75$ 的样本量下经 McNemar 精确检验均\emph{未达统计显著}（$p\ge0.25$）；因此我们把它们定位为“低开销、零回归的确定性工程杠杆”，而非已被统计确立的性能提升。

\subsection{F3: 跨应用泛化与 BEM 盲区}
\label{sec:f3}

\subsubsection{设置}
任何仅在单一应用上完成的试点都可能被审稿人合理质疑：所学到的可能只是应用专有的约定。为此本文以 \texttt{git sparse-checkout} 方式仅检出 \texttt{payloadcms/payload}~\cite{payloadCMS} 的最新版本（HEAD；ReproBreak 中第二大的 Playwright 案例组，包含 319 个失效配对），在 \texttt{packages/ui/src} 上运行未修改的 L1 抽取器，产出 457 条记录。对每个失效配对，使用 \texttt{new\_locator} 字符串在该记录表上做解析匹配（以 HEAD 版本可达率作为替代评估指标，逐 commit 配对未纳入本文）。

\subsubsection{结果}
\begin{table}[t]
\centering
\caption{跨应用 L1 可达率（\S\ref{sec:f3}）。原始下降是真实的，但成因可诊断。}
\label{tab:cross-app}
\small
\resizebox{\columnwidth}{!}{%
\begin{tabular}{lrr}
\toprule
指标 & \texttt{koenig} & \texttt{payload} \\
\midrule
抽取器记录数（均值/单次） & 每 commit 约 400 & 457（HEAD） \\
Playwright 失效配对 & 93（已复现） & 319（CSV） \\
\texttt{new\_locator} 可达 & 75 / 93 = \textbf{80.6\%} & 12 / 319 = \textbf{3.8\%} \\
\midrule
主导锚点 & \texttt{data-kg-*}（约 50\%） & CSS class$^\dagger$ \\
模板字面量 class? & 罕见 & 普遍（BEM） \\
\bottomrule
\end{tabular}}
\\[2pt]
{\scriptsize $^\dagger$payload 的主导锚点占比来自对 307 个未命中用例的采样，而非全量失效分布。}
\end{table}

原始 3.8\% 的可达率很低，但对 307 个未命中用例的检查显示，其原因十分明确。payload 使用 BEM~\cite{bemMethodology}: 每个组件声明一个 \texttt{baseClass} 常量（例如 \texttt{const baseClass = 'dashboard'}），并以 \texttt{\{`\$\{baseClass\}\_\_add-button`\}} 这样的模板字面量写出 className。测试引用\emph{渲染后}的类 \texttt{.dashboard\_\_add-button}，这对一个纯字面量的 AST 抽取器是不可见的。

\subsubsection{解读}
这一结果并非对 HealReact L1 设计的否证，而是识别出一处范围明确的盲区，其修复对应一次确定性的 AST 处理：构造一个 \texttt{baseClass} 解析器，回溯到常量声明并展开其模板字面量，并配以一个 CSS Modules~\cite{cssModules} 等价的展开器。基于对 307 个未命中用例的\emph{非系统}采样，本文\emph{假设}修复后原始可达率会显著抬升，但该结论尚属假设而非测量；对全部 307 个未命中用例做系统化的失因标注（BEM、CSS Modules、运行时 DOM、解析器局限等类别计数）、固定并记录 payload 的 HEAD commit SHA 以及与 ReproBreak 的逐 commit 配对对齐，均留作后续工作（当前冻结工件仅记录了 payload 的仓库、配对数与记录数，未记录精确 HEAD SHA，属已知复现缺口）。3.8\% 是当前抽取器下诚实的原始测量值。该发现真正的贡献是这条可证伪的跨应用泛化\emph{假设}本身：锚点文化的多样性（\texttt{data-testid} vs 自定义 \texttt{data-*} vs BEM 模板字面量）很可能是定位器修复工具泛化能力的关卡，任何仅在单一应用上完成基准测试的工具都可能高估其泛化性。
